\documentclass[10pt,a4paper]{amsart}
\usepackage[margin=19mm]{geometry}
\usepackage{amsmath,amssymb,mathtools}
\usepackage[scr=rsfs,cal=euler]{mathalfa}
\usepackage{xfrac}
\usepackage[unicode,hidelinks,bookmarks=false]{hyperref}
\newcommand{\ie}{i.e.\ }
\newcommand{\cf}{cf.\ }
\newcommand{\resp}{respectively\ }
\newcommand{\Subsection}{\subsection}
\providecommand{\nobreakdash}{\nobreak\hbox{-}}
\setlength{\emergencystretch}{2em}
\multlinegap0pt
\usepackage{bm}
\usepackage{bbm}
\usepackage{dsfont}
\usepackage{xspace}
\usepackage{cancel}
\usepackage{mathtools}
\usepackage{amsthm}
\makeatletter
\@ifundefined{lemma}{\newtheorem{lemma}{Lemma}}{}
\@ifundefined{proposition}{\newtheorem{proposition}{Proposition}}{}
\@ifundefined{theorem}{\newtheorem{theorem}{Theorem}}{}
\makeatother
\title[Central $L$ values of congruent number elliptic curves]{Central $L$ values of congruent number elliptic curves}
\author{Xuejun Guo, Dongxi Ye, Hongbo Yin}
\date{Source: arXiv:2505.13133v2 (CC BY 4.0)}
\begin{document}
\maketitle

\noindent\emph{Source and licence.} Central $L$ values of congruent number elliptic curves. Version arXiv:2505.13133v2 (CC BY 4.0), distributed under the Creative Commons Attribution 4.0 International licence, as recorded in the arXiv licence metadata for this exact version and shown on the abstract page https://arxiv.org/abs/2505.13133v2. Full rights notice: licenses/arxiv-2505.13133-CC-BY-4.0.txt.\par\smallskip

\noindent\textit{Editorial scope.} Counted unit: the Theorem whose source label is \texttt{main4} in the article (source0.tex lines 1874--1889, source label \texttt{main4}), delivered together with the article's own complete original proof of it (source0.tex lines 1891--1941, one \texttt{proof} environment of 51 lines). No mathematics is written, repaired or re-derived: statement and proof are the article's own text line for line. Required context retained uncounted: main2 (lines 1091--1113); corthetaf (lines 1670--1693); lemna (lines 1755--1761); shimuraff (lines 1814--1844). Every definition, display and label of those lines is retained unchanged. Omitted from this package and not redistributed: 1--1090, 1114--1669, 1694--1754, 1762--1813, 1845--1873, 1890--1890, 1942--2467. Nothing inside a retained excerpt is omitted or altered: every retained body is delivered exactly as the article's lines are written, so no omission receipt of any kind accompanies this package. Citations: the retained excerpts carry no citation command at all, so no bibliography entry of the article is transcribed and no reference is redistributed. Reference adaptation: none was needed and none was made; every reference inside the delivered unit resolves inside it, so no unresolved reference or citation is produced. Printed slips kept literal: no printed word, symbol, hypothesis or number of the counted statement or of its proof was altered. Layout and class: the article's own class is \texttt{amsart} with packages including amssymb, stmaryrd, epsfig, mathdots, hyperref, fullpage, float, color, diagbox, slashbox, tikz, booktabs, tikz-cd, pgfplots; the wrapper is the lane's pinned \texttt{amsart} editorial wrapper at 10pt on A4, which restates the theorem-like environments the retained text needs and adds the article's own macro definitions that the retained text needs (none), with extra wrapper packages (bm, bbm, dsfont, xspace, cancel, mathtools). The delivered numbering is the unit's own. Assets: no retained span contains a figure environment, a graphics-inclusion command, a PDF-page inclusion, a table or a TikZ picture; no image file is copied into this package, the package declares no asset dependency and no image file is redistributed with it. Licence: version arXiv:2505.13133v2 of this article is distributed under the Creative Commons Attribution 4.0 International licence, as recorded in the arXiv licence metadata for this exact version and shown on the abstract page https://arxiv.org/abs/2505.13133v2. A full rights notice is delivered at licenses/arxiv-2505.13133-CC-BY-4.0.txt. Characters: every retained excerpt is pure ASCII, so the delivered engine, XeLaTeX with its default Latin Modern text font, reports no missing character.
\begin{theorem}\label{main2}
Write 
$$
n'=\begin{cases}
    n,&\mbox{if $n$ is odd};\\
    m,&\mbox{if $n=2m$.}
\end{cases}
$$
and denote by ${\rm rad}(n')$ the radical of~$n'$. Then
$$
L(E_{n},1)=\begin{cases}
    \displaystyle{\frac{\pi}{4}\Theta((1+i)/2)\frac{1}{\sqrt{n}}{\rm Tr}_{H_{\mathfrak{m}}/K}(\sqrt{n}\Theta_{n}(1+i))},&\mbox{if $n$ is odd};\\
    \displaystyle{\frac{\pi}{4}\Theta(i/2)\frac{1}{\sqrt{n}}{\rm Tr}_{H_{\mathfrak{m}}/K}(\sqrt{n}\Theta_{n}(i))},&\mbox{if $n=2m$},
\end{cases}
$$
   %     $$
%     \frac{L(1,\tilde{\chi}_{n})}{\frac{\pi}{4}(\Theta(1+i)+\Theta_{\chi_{-4}}((1+i)/2))}=\frac{1}{\sqrt{n}}{\rm Tr}_{H_{\mathfrak{m}}/K}(\sqrt{n}\Theta_{n}(1+i)),
%     $$
    where
    $$
    \Theta_{n}(\tau)=\sum_{\substack{p_{1}\cdots p_{\ell}|{\rm rad}(n')\\ p_{1}\cdots p_{\ell}\equiv r\pmod{4}}}\frac{(-1)^{\ell+\frac{r-1}{2}}}{p_{1}\cdots p_{\ell}}\frac{\Theta\left(\frac{n'}{p_{1}\cdots p_{\ell}}\tau/2\right)}{\Theta(\tau/2)}.
    $$
\end{theorem}

\begin{proposition}\label{corthetaf}
      Write 
$$
n'=\begin{cases}
    n,&\mbox{if $n$ is odd};\\
    m,&\mbox{if $n=2m$}.
\end{cases}
$$   
and denote by ${\rm rad}(n')$ the radical of~$n$, and let 
    $$
    \mathcal{A}=[a,b+i]\quad\mbox{and}\quad \mathcal{A}_{1}=[a_{1},b+i]
    $$
    be two primitive integral ideals with $b$ being  such that $b^{2}\equiv-1\pmod{n'a^{2}a_{1}}$. Then
    \begin{align*}
       % & \sum_{p_{1}\cdots p_{\ell}|{\rm rad}(n')}\frac{(-1)^{\ell}}{p_{1}\cdots p_{\ell}}\frac{\Theta\left(\frac{n'}{p_{1}\cdots p_{n}}\tau_{\mathcal{A}}/2\right)}{\Theta(\tau_{\mathcal{A}}/2)}\\
       % &= \frac{1}{n}\sum_{p_{1}\cdots p_{\ell}|{\rm rad}(n')}(-1)^{\ell}\sum_{r=0}^{\frac{n}{p_{1}\cdots p_{n}}-1}f_{ar,\frac{n}{p_{1}\cdots p_{\ell}}}(\tau_{\mathcal{A}^{2}\mathcal{A}_{1}})\overline{f_{r,\frac{n}{p_{1}\cdots p_{\ell}}}(\tau_{\mathcal{A}_{1}})}\\
       % &=\frac{1}{n}\sum_{p_{1}\cdots p_{\ell}|{\rm rad}(n')}(-1)^{\ell}\sum_{r=0}^{\frac{n}{p_{1}\cdots p_{n}}-1}f_{a(p_{1}\cdots p_{\ell}r),n}(\tau_{\mathcal{A}^{2}\mathcal{A}_{1}})\overline{f_{(p_{1}\cdots p_{\ell}r),n}(\tau_{\mathcal{A}_{1}})}\\
       \Theta_{n}(\tau_{\mathcal{A}})&=\frac{1}{n'}\sum_{p_{1}\cdots p_{\ell}|{\rm rad}(n')}(-1)^{\ell}+\frac{1}{n'}\sum_{r\in(\mathbb{Z}/n'\mathbb{Z})^{\times}}f_{ar,n'}(\tau_{\mathcal{A}^{2}\mathcal{A}_{1}})\overline{f_{r,n'}(\tau_{\mathcal{A}_{1}})},
    \end{align*}
    where $\Theta_{n}(\tau)$ is defined as in Theorem~\ref{main2}, and
    $$
    f_{r,d}(\tau)=\frac{\theta_{[r/d,1/2]}(\tau)}{\theta_{[0,1/2]}(\tau)}.
    $$
\end{proposition}

\begin{lemma}\label{lemna}
Given a positive odd integer~$n$.    For a class $[\mathcal{A}]\in J_{\mathfrak{m}}^{*}/J_{\mathfrak{m},\mathbb{Z}}^{*}$,  there is a representative $\mathcal{A}$ such that $\mathcal{A}=\alpha_{\mathcal{A}}\mathcal{O}_{K}=[a,b+i]$, and $n_{a}\equiv1\pmod{n}$, where $n_{a}$ is defined by
$$
n_{a}a+m_{a}(b+i)=\alpha_{\mathcal{A}}.
$$

\end{lemma}

\begin{lemma}\label{shimuraff}
% For $[\mathcal{A}]\in J_{4n}^{*}/J_{\mathfrak{m},\mathbb{Z}}^{*}$ with $\mathcal{A}=\alpha_{\mathcal{A}}=[a,b+i]\equiv1\pmod{4}$ primitive, where $b$ is even and such that $b^{2}=-1\pmod{n^{2}a}$, one has
For a positive odd integer~$D|n$, let $\alpha\mathcal{O}_{K}=[a,b+i]=[a,a\tau_{\mathcal{A}}]$ be a primitive integral ideal with $\alpha\equiv1\pmod{4}$,  $a\equiv1\pmod{4}$, $b$  such that $b^{2}\equiv-1\pmod{D^{2}a}$, and $n_{a}\equiv1\pmod{4D}$ given by Lemma~\ref{lemna}. Then one has
$$
f_{r,D}(\tau_{\mathcal{A}})=f_{r,D}(b+i)
$$
for any~$r$.
% $$
% f_{r}(\tau_{\mathcal{A}})=e^{2\pi i\left(\frac{(t_a-1)n_{a}2r}{4p}\right)}f_{n_{a}r}(b+i)=e^{2\pi i\left(\frac{(1-n_{a})2r}{4p}\right)}f_{n_{a}r}(b+i)=f_{r}(b+i)^{\sigma_{\mathcal{A}}},
% $$
% where $t_{a}$ and $n_{a}$ are defined by
% $$
% t_{a}a+s(b+i)=\alpha_{\mathcal{A}},\quad n_{a}=at_a+2bs\pmod{p}\,\,\mbox{chosen to be odd}
% $$
% which are induced by
% $$
% \alpha_{\mathcal{A}}\begin{pmatrix}
%     b+i\\1
% \end{pmatrix}=\begin{pmatrix}
%     t_{a}a+2bs&-s({b^{2}+1})\\
%     s&t_{a}a
% \end{pmatrix}\begin{pmatrix}
%     b+i\\1
% \end{pmatrix}.
% $$
% So $(t_{a}a+2bs)t_{a}=n_{a}t_{a}\equiv 1\pmod{p^{2}}$.
% One therefore has 
% $$
% e^{-2\pi i\frac{2r}{4p}}f_{r}(\tau_{\mathcal{A}})=e^{-2\pi i\frac{2n_{a}r}{4p}}f_{n_{a}r}(b+i).
% $$
\end{lemma}

\begin{theorem}\label{main4}
\begin{enumerate}
    \item For~$n$ being odd, and an even number~$b$ such that $b^{2}=-1\pmod{n^{2}}$, one has
    \[
L(E_n,1)=\frac{\pi \left|\theta_{\chi_{n}}\left(\frac{b+i}{2
n^2}\right)\right|^2}{4\sqrt{2}n }.
\]

\item For~$n=2m$, and an odd number~$b$ such that $b^{2}=-1\pmod{m^{2}}$, one has
    \[
L(E_n,1)=\frac{\pi \left|\theta_{\chi_{m}}\left(\frac{b+i}{2
m^2}\right)\right|^2}{\sqrt{2}n }.
\]
\end{enumerate}

\end{theorem}

\begin{proof}
%     Given Lemma~\ref{thetathetagalois}, one has
% \begin{align*}
% &\frac{1}{\sqrt{n}}{\rm Tr}_{H_{\mathfrak{m}}/K}\left(\sqrt{n}\Theta_{n}(1+i)\right)
% % &=\sum_{\substack{[\mathcal{A}]\in J_{\mathfrak{m}}^{*}/J_{\mathfrak{m},\mathbb{Z}}^{*}\\\mathcal{A}=[a,b+i]\\b\,\,odd}}\chi'_{n}({\mathcal{A}})\sum_{p_{1}\cdots p_{\ell}|{\rm rad}(n')}\frac{(-1)^{\ell}}{p_{1}\cdots p_{\ell}}\left(\frac{\Theta\left(\frac{n}{p_{1}\cdots p_{\ell}}\tau_{\mathcal{A}}\right)+\Theta_{\chi_{-4}}\left(\frac{n}{p_{1}\cdots p_{\ell}}\tau_{\mathcal{A}}/2\right)}{\Theta(\tau_{\mathcal{A}})+\Theta_{\chi_{-4}}(\tau_{\mathcal{A}}/2)}\right)\\
%     =\sum_{\substack{[\mathcal{A}]\in J_{\mathfrak{m}}^{*}/J_{\mathfrak{m},\mathbb{Z}}^{*}\\ \mathcal{A}=[a,b+i]\\b\,\,odd\\b^{2}\equiv-1\pmod{n^{2}a^{2}}\\n_{a}\equiv1\pmod{4n}}}\chi_{n}'(\mathcal{A})\Theta_{n}(\tau_{\mathcal{A}}).
% \end{align*}

    
We only prove part (1) as an illustration.  Note that the proof of Theorem~\ref{main2} together with Lemma~\ref{lemna} yields that
     $$
   L(E_{n},1)=\frac{\pi\Theta((1+i)/2)}{4}\sum_{\substack{[\mathcal{A}]\in J_{\mathfrak{m}}^{*}/J_{\mathfrak{m},\mathbb{Z}}^{*}\\ \mathcal{A}=[a,b+i]\\b\,\,odd\\b^{2}\equiv-1\pmod{n^{2}a^{2}}\\n_{a}\equiv1\pmod{4n}}}\chi_{n}'(\mathcal{A})\Theta_{n}(\tau_{\mathcal{A}}).
    $$
Upon this,    using Proposition~\ref{corthetaf} and the simple fact of
$$
\sum_{[\mathcal{A}]\in J_{\mathfrak{m}}^{*}/J_{\mathfrak{m},\mathbb{Z}}^{*}}\chi'_{n}(\mathcal{A})=0,
$$
as well as Lemma~\ref{shimuraff},  one can deduce from the equality above that
   \begin{align*}
        L(E_{n},1)&=\frac{\pi\Theta((1+i)/2)}{4n}\sum_{\substack{[\mathcal{A}]\in J_{\mathfrak{m}}^{*}/J_{\mathfrak{m},\mathbb{Z}}^{*}\\ \mathcal{A}=[a,b+i]\\b\,\,odd\\b^{2}\equiv-1\pmod{n^{2}a^{2}}\\n_{a}\equiv1\pmod{4n}}}\sum_{r\in(\mathbb{Z}/n\mathbb{Z})^{\times}}\chi_{n}'(\mathcal{A})f_{ar,n}(b+i)\overline{f_{r,n}(b+i)}\\
         &=\frac{\pi\Theta((1+i)/2)}{4n}\sum_{a\in(\mathbb{Z}/n\mathbb{Z})^{\times}}\sum_{r\in(\mathbb{Z}/n\mathbb{Z})^{\times}}\chi_{n}(ar)f_{ar,n}(b+i)\overline{\chi_{n}(r)f_{r,n}(b+i)}\\
         &=\frac{\pi\Theta((1+i)/2)}{4n}\left|\sum_{r=1}^{n-1}\chi_{n}(r)f_{r,n}(\tilde{b}+i)\right|^{2}
        % &=\frac{\left|\theta_{\chi_{n}}\left(\frac{b+i}{2n^{2}}\right)\right|^{2}}{n\theta_{[0,1/2]}\left(b+i\right)^{2}}\\
         %&=\frac{\left|\theta_{\chi_{n}}\left(\frac{b+i}{2n^{2}}\right)\right|^{2}}{n\theta\left(i/2\right)^{2}}
         %        \frac{1}{p}\sum_{s=1}^{p-1}\sum_{r=1}^{p-1}\chi_{n}(s)f_{sr}(b+i)\overline{f_{r}(b+i)}.
    \end{align*}
    for a $\tilde{b}$ being odd such that $\tilde{b}^{2}+1=0\pmod{n^{2}}$. Finally, by the definition of $f_{r,n}(\tau)$, it is routine to check that
    $$
    \sum_{r=1}^{n-1}\chi_{n}(r)f_{r,n}(\tilde{b}+i)=\frac{\theta_{\chi_{n}}\left(\frac{b+i}{2n^{2}}\right)}{\theta(i/2)}
    $$
    for a $b$ being even such that $b^{2}+1=0\pmod{n^{2}}$. Finally, noting that
    $$
    \Theta\left(\frac{1+i}{2}\right)=\theta\left(\frac{1+i}{2}\right)^{2}
    $$
    and making use of the classical result
    $$
    \frac{\theta(i/2)^{2}}{\theta((1+i)/2)^{2}}=\sqrt{2}
    $$
    give the desired formula.

%      For part (2),  one similarly has that
%      $$
%    L(E_{n},1)=\frac{\pi\Theta(i/2)}{4m}\frac{\left|\theta_{\chi_{n}}\left(\frac{b+i}{2m^{2}}\right)\right|^{2}}{\theta((1+i)/2)^{2}}.
%     $$
% Again, using the classical result
%     $$
%     \frac{\theta((1+i)/2)^{2}}{\theta(i/2)^{2}}=\frac{\sqrt{2}}{2}
%     $$
%     and recalling that $2m=n$
%     give the desired formula.
\end{proof}

\end{document}
